\section{Introduction}\label{sec:intro}

\subsection{Motivation}\label{sec:motivation}

Search and rescue (SAR) operations in wilderness and mountain environments are critically time-sensitive. The \textit{golden hour}---the window within which medical intervention most significantly improves survival outcomes---places severe pressure on response teams to locate casualties rapidly \cite{goodrich2008}. In the United Kingdom, mountain rescue teams responded to over 3{,}000 incidents in 2023, with search phases often lasting several hours across terrain that is difficult or dangerous to traverse on foot \cite{mra2023}. During this time, casualties may be exposed to hypothermia, haemorrhage, or other deteriorating conditions that reduce the probability of survival.

The challenge is fundamentally one of coverage rate. A ground search team of ten members can sweep approximately $0.1$--$0.5$\,km$^{2}$/hr depending on terrain complexity and vegetation density, whereas a single UAV at 30\,m altitude with a \ang{60} field of view can survey a 35\,m-wide swath at 5\,m/s, covering approximately $0.6$\,km$^{2}$/hr---an order-of-magnitude improvement per searcher deployed \cite{goodrich2008, murphy2014}. This advantage is particularly pronounced in the first hours of an incident, when the search area is at its largest and the probability of survival is highest. UK SAR statistics indicate that the median time from alert to casualty contact exceeds 90\,minutes in mountain terrain, a figure that has remained broadly static despite improvements in team coordination and training \cite{mra2023}. Reducing this interval through faster aerial survey represents a direct and measurable improvement in outcomes.

Unmanned aerial vehicles (UAVs) offer a compelling solution. A drone equipped with a camera and onboard computer vision can survey large areas systematically at altitude, unimpeded by terrain \cite{murphy2014}. Recent advances in lightweight object detection models, particularly the YOLO family \cite{yolohistory, yolov8}, have made it feasible to run real-time inference on edge devices such as the Raspberry Pi, enabling fully onboard detection without dependence on a ground-station data link \cite{edgeai}.

\subsection{Related Work}\label{sec:related-work}

\subsubsection{UAVs in Search and Rescue}

The use of UAVs for SAR has been studied extensively over the past two decades. Early work by Goodrich et al.\ \cite{goodrich2008} demonstrated a camera-equipped mini-UAV supporting wilderness search operations, but relied entirely on a human operator to visually scan a live video feed for casualties. Murphy's comprehensive survey of disaster robotics \cite{murphy2014} catalogues numerous deployments of UAVs in post-disaster scenarios, noting that the primary role of UAVs has been to provide situational awareness rather than autonomous detection. Silvagni et al.\ \cite{sarreview} developed a multipurpose UAV for avalanche SAR using thermal imaging, achieving detection of buried casualties, but at a hardware cost that limits widespread adoption by volunteer mountain rescue teams.

Commercial SAR drone platforms have matured considerably. The DJI Matrice series, equipped with thermal and RGB payloads, is now standard equipment for many professional SAR teams \cite{dji_matrice}, and SenseFly's fixed-wing mapping drones have been deployed for large-area survey following natural disasters \cite{sensefly_sar}. However, these platforms remain expensive (typically \$5{,}000--\$30{,}000), require trained operators for both piloting and image interpretation, and depend on high-bandwidth data links for real-time imagery review. Mishra et al.\ \cite{mishra2020} survey drone applications in disaster management, concluding that while UAV technology is mature for reconnaissance and mapping, autonomous casualty detection remains an open challenge.

Thermal imaging is the most widely studied sensing modality for aerial SAR \cite{thermal_sar, thermal_review}. Burke et al.\ \cite{thermal_sar} optimised UAV flight parameters for thermal detection, finding that altitude, time of day, and ground cover significantly affect reliability. Al-Kaff et al.\ \cite{alkaff2019} combined thermal and RGB cameras for person detection from UAVs, achieving improved performance through sensor fusion. While thermal cameras offer the advantage of detecting body heat regardless of visual camouflage, they add \$2{,}000--\$10{,}000 to the platform cost, have lower spatial resolution than RGB cameras at equivalent price points, and suffer from reduced effectiveness when ambient temperatures approach body temperature \cite{thermal_review}. This project deliberately avoids thermal sensing in favour of RGB-only detection to explore the feasibility of a lower-cost approach.

\subsubsection{Object Detection from Aerial Platforms}

The YOLO (You Only Look Once) family of detectors has become the dominant architecture for real-time aerial object detection \cite{yolohistory}. The original YOLO introduced the single-pass detection paradigm, processing an entire image in one forward pass through a convolutional neural network and achieving frame-rate inference on GPU hardware \cite{yolohistory}. Subsequent iterations---YOLOv3 \cite{yolov3}, YOLOv5, and YOLOv8 \cite{yolov8}---have progressively improved accuracy while reducing model size, with YOLOv8n (the \textit{nano} variant) requiring only 3.2\,million parameters and 8.7\,GFLOPs.

Aerial person detection presents distinct challenges compared to ground-level detection. Subjects appear at much smaller pixel scales, aspect ratios differ from the upright pedestrians typical of ground-level datasets such as COCO \cite{coco_dataset}, and background clutter from natural terrain introduces frequent false positives \cite{zhu2021}. Zhu et al.\ \cite{zhu2021} survey object detection in aerial images, identifying small object scale, orientation variation, and class imbalance as the three principal challenges. Mandal et al.\ \cite{mandal2020} demonstrated person detection from a drone at 10--50\,m altitude using a fine-tuned SSD model, achieving 72\% mAP but noting a sharp decline above 30\,m. The VisDrone dataset and benchmark \cite{visdrone2021} have established a standard evaluation framework for aerial object detection, confirming that even state-of-the-art detectors achieve significantly lower mAP on aerial imagery than on ground-level benchmarks.

\subsubsection{Edge AI for UAV Applications}

Running neural network inference onboard a UAV, rather than streaming video to a ground station for processing, eliminates the dependency on a continuous high-bandwidth data link and reduces latency from detection to action. This is the domain of \textit{edge AI}---deploying trained models on resource-constrained devices at the point of data acquisition \cite{edgeai}.

The NVIDIA Jetson family of embedded GPU platforms has been the most widely studied hardware for onboard UAV inference. Tijtgat et al.\ \cite{tijtgat2017} demonstrated real-time person detection on a Jetson TX2 mounted on a drone, achieving 10--15\,FPS with a compressed YOLO model. Anwar and Raychowdhury \cite{anwar2020} surveyed autonomous aerial navigation with deep learning, finding that Jetson-class hardware could sustain real-time inference for obstacle avoidance and object detection, but at a power draw of 10--15\,W and a unit cost of \$200--\$400.

An alternative approach uses quantised models on lower-power hardware. TensorFlow Lite \cite{tflite} enables deployment of neural networks on ARM processors with minimal dependencies, using XNNPACK \cite{xnnpack} for optimised floating-point inference. The Raspberry Pi~5, with its quad-core Cortex-A76 processor, represents the low-cost extreme at approximately \pounds{}70, consuming under 5\,W \cite{raspberrypi5}. While inference speed is substantially lower than on GPU-accelerated platforms---typically 4--5\,FPS for YOLOv8n on a Pi~5 versus 30+\,FPS on a Jetson Nano---this is acceptable for a search mission where the drone traverses slowly and each ground patch remains in the camera's field of view for several seconds. Mazzia et al.\ \cite{mazzia2022} survey edge AI deployment strategies, concluding that model quantisation and hardware-aware architecture search are the key enablers for sub-watt inference on ARM devices.

\subsubsection{Coverage Path Planning}

Systematic area coverage is a fundamental requirement for autonomous search. The boustrophedon (lawnmower) pattern---parallel sweeps with \ang{180} turns at the boundary---is the canonical coverage strategy and has been shown to be optimal for convex regions under uniform detection probability \cite{boustrophedon}. Galceran and Carreras \cite{covpath} survey coverage path planning comprehensively, classifying approaches into exact cellular decomposition, approximate decomposition, and graph-based methods. For convex or near-convex search areas with uniform terrain, boustrophedon decomposition provides complete coverage with minimal redundancy.

In the SAR context, sweep width must be matched to the sensor footprint at the planned altitude to ensure no gap exists between adjacent passes. Hayat et al.\ \cite{multiuav} survey UAV network architectures for civil applications including SAR, noting that multi-UAV coordination can reduce search time linearly with fleet size but introduces communication and deconfliction complexity. For a single-UAV system such as the one presented here, the planning problem reduces to computing a lawnmower pattern that covers a user-defined polygon with appropriate lane spacing.

\subsection{Research Gap and Contributions}\label{sec:contributions}

Despite the substantial body of work on UAVs for SAR, a gap persists between the capability demonstrated in research and the systems available to operational SAR teams. The majority of published systems fall into one of three categories: (1)~platforms that stream video to a ground station for human interpretation, requiring a skilled operator and high-bandwidth link \cite{goodrich2008}; (2)~systems that employ thermal cameras, adding significant cost \cite{thermal_sar, sarreview}; or (3)~research prototypes using GPU-class onboard computers (Jetson TX2/Xavier) that are too heavy, power-hungry, or expensive for adoption by volunteer rescue organisations \cite{tijtgat2017, anwar2020}.

What is absent is a demonstration that a \textit{low-cost, fully autonomous} SAR drone can be built using commodity hardware---specifically, a Raspberry Pi as the onboard computer and an RGB camera as the sole sensor---with all detection and decision-making performed onboard without cloud or ground-station processing. This project addresses that gap directly.

The specific contributions of this work are:

\begin{enumerate}
    \item \textbf{End-to-end autonomous SAR system on commodity hardware.} A complete mission pipeline---from takeoff through coverage search, onboard AI-based casualty detection, target approach, operator verification, and landing---running on a Raspberry Pi~5 with a total onboard computing cost under \pounds{}100.

    \item \textbf{Custom-trained YOLOv8n model for aerial casualty detection.} A lightweight object detection model trained on a combination of synthetic and real aerial imagery at native camera resolution ($1456 \times 1088$), achieving 0.995 mAP$_{50}$ while running at 4.8\,FPS on the Pi~5 via TensorFlow Lite.

    \item \textbf{GPS-based target localisation from visual detections.} A method for estimating the ground-level GPS coordinates of a detected casualty by combining the drone's GPS position, altitude, and attitude with the pixel location of the detection and calibrated camera intrinsics, using spatial clustering and Kalman filtering to refine estimates across multiple observations.

    \item \textbf{Progressive testing methodology for autonomous UAV systems.} A five-stage testing pipeline---from software-in-the-loop simulation through bench testing, passive flight observation, autonomous flight with logging only, to full autonomous operation---that systematically verifies each subsystem before integration, reducing the risk of failure during live flight.

    \item \textbf{Human-in-the-loop verification architecture.} A web-based ground station that presents detection imagery to a remote operator for confirmation before the drone commits to landing, balancing autonomy with operational safety and addressing the false-positive problem inherent in RGB-only detection.
\end{enumerate}

\subsection{Objectives}\label{sec:objectives}

The primary objective of this project is to demonstrate an autonomous SAR mission in which a UAV:

\begin{enumerate}
    \item Takes off and flies a lawnmower search pattern over a defined area;
    \item Uses an onboard YOLOv8-based object detection model to identify a casualty (represented by a life-size dummy);
    \item Centres on the detected target, descends to a verification altitude, and presents the detection to a remote operator for confirmation;
    \item Upon operator confirmation, approaches and lands at a safe offset from the target.
\end{enumerate}

A human-in-the-loop verification step ensures that the drone autonomously detects and approaches, but a human operator makes the final confirmation decision before landing, balancing autonomy with safety.

A secondary objective is to validate the system through a progressive testing methodology, advancing from software-in-the-loop (SITL) simulation through bench testing with real hardware to outdoor flight trials, ensuring that each subsystem is verified independently before integration.

\subsection{Scope and Constraints}\label{sec:scope}

This work is conducted as part of the AENGM0074 module at the University of Bristol, undertaken by a team of five MSc students over a single academic term. Several constraints shape the design:

\begin{itemize}
    \item \textbf{Budget:} The project uses a Raspberry Pi~5 as the onboard computer and a Pixhawk Cube autopilot, both selected for cost-effectiveness. No thermal imaging or high-end GPU hardware is employed.
    \item \textbf{Regulatory:} All flights comply with UK Civil Aviation Authority (CAA) regulations. The demonstration flight takes place within a designated university-approved flight area with appropriate risk assessments in place \cite{ukcaa}.
    \item \textbf{Detection target:} A life-size dummy is used as the casualty surrogate, enabling repeatable testing. The custom YOLOv8n model is trained on a combination of synthetic and real imagery at native camera resolution.
    \item \textbf{Single demonstration:} Due to time and weather constraints, the system is designed around a single outdoor flight demonstration, preceded by extensive simulation and bench testing.
    \item \textbf{Communication:} The operator interface is a web-based ground station accessed over a local Wi-Fi link, with no dependence on cellular or internet connectivity.
\end{itemize}

The system is a proof-of-concept prototype, not intended for deployment in operational SAR, but rather to demonstrate the feasibility of the approach and to identify the key engineering challenges---detection reliability at altitude, GPS estimation accuracy, and inference speed on edge hardware---that must be addressed in a production system.

\subsection{Report Structure}\label{sec:report-structure}

The remainder of this report is organised as follows. Section~\ref{sec:architecture} presents the system architecture and software modularity. Section~\ref{sec:hardware} describes the hardware platform: the Raspberry Pi~5 companion computer, the Cube Orange+ autopilot, and the IMX296 global shutter camera. Section~\ref{sec:cv} details the computer vision pipeline, covering model selection, the dual-backend inference architecture, and preprocessing including lens distortion correction. Section~\ref{sec:planning} describes the lawnmower coverage path planning algorithm for arbitrary polygonal search areas. Section~\ref{sec:statemachine} presents the finite state machine governing the mission. Section~\ref{sec:localisation} describes GPS-based target localisation via pixel-to-GPS projection, spatial clustering, and Kalman filtering. Section~\ref{sec:groundstation} covers the web-based ground station and operator interface. Section~\ref{sec:simulation} describes the software-in-the-loop simulation environment. Section~\ref{sec:testing} presents the progressive testing methodology. Section~\ref{sec:results} analyses field trial results. Section~\ref{sec:training} details the model training pipeline, including synthetic data generation. Section~\ref{sec:safety} discusses safety and risk management. Section~\ref{sec:future} identifies directions for future work, and Section~\ref{sec:conclusion} draws conclusions.
